\documentclass[11pt,a4paper]{article}
% Wspólny szablon dla kart supermocy (jedna karta = jedna strona A4)
% Użycie w karcie: \documentclass[11pt,a4paper]{article} % Wspólny szablon dla kart supermocy (jedna karta = jedna strona A4)
% Użycie w karcie: \documentclass[11pt,a4paper]{article} % Wspólny szablon dla kart supermocy (jedna karta = jedna strona A4)
% Użycie w karcie: \documentclass[11pt,a4paper]{article} \input{../preambula} \begin{document}

% Silnik tectonic (XeTeX): polskie znaki przez fontspec, nie przez inputenc/T1
\usepackage{fontspec}
\setmainfont{Arial}
\usepackage[polish]{babel}
\usepackage[a4paper,margin=1.5cm]{geometry}
\usepackage{amsmath,amssymb}
\usepackage{tikz}
\usetikzlibrary{arrows.meta,positioning,calc,shapes.geometric,decorations.pathreplacing}
\usepackage[most]{tcolorbox}
\usepackage{enumitem}
\usepackage{booktabs}
\usepackage{xcolor}

\pagestyle{empty}
\setlength{\parindent}{0pt}
\setlength{\parskip}{0.4em}

\definecolor{cel}{HTML}{1F4E79}
\definecolor{regula}{HTML}{E8F1FA}
\definecolor{pulapka}{HTML}{FDECEA}
\definecolor{okres}{HTML}{C0392B}
\definecolor{skonczony}{HTML}{1E8449}

% Ramka na regułę
\newtcolorbox{regula}[1][]{
  colback=regula, colframe=cel, boxrule=0.8pt, arc=2pt,
  left=6pt, right=6pt, top=4pt, bottom=4pt, title=#1,
  fonttitle=\bfseries, coltitle=white
}

% Ramka na pułapkę z kartkówki
\newtcolorbox{pulapka}[1][]{
  colback=pulapka, colframe=okres, boxrule=0.8pt, arc=2pt,
  left=6pt, right=6pt, top=4pt, bottom=4pt, title=#1,
  fonttitle=\bfseries, coltitle=white
}

% Nagłówek karty: numer, tytuł, cel
\newcommand{\naglowekKarty}[3]{%
  {\LARGE\bfseries\color{cel} Karta #1: #2}\\[2pt]
  {\large Cel: #3}\par
  \vspace{0.3em}\noindent\rule{\linewidth}{1pt}\par
}

% Pole na obliczenia ucznia: ramka z delikatną kratką, argument = wysokość
\newcommand{\miejsce}[1]{%
  \par\vspace{0.2em}\noindent
  \begin{tikzpicture}
    \draw[step=0.5cm, gray!25, very thin] (0,0) grid (\textwidth,#1);
    \draw[gray!60, rounded corners=2pt] (0,0) rectangle (\textwidth,#1);
    \node[anchor=north west, font=\scriptsize, text=gray] at (0.05,#1) {obliczenia};
  \end{tikzpicture}\par}
 \begin{document}

% Silnik tectonic (XeTeX): polskie znaki przez fontspec, nie przez inputenc/T1
\usepackage{fontspec}
\setmainfont{Arial}
\usepackage[polish]{babel}
\usepackage[a4paper,margin=1.5cm]{geometry}
\usepackage{amsmath,amssymb}
\usepackage{tikz}
\usetikzlibrary{arrows.meta,positioning,calc,shapes.geometric,decorations.pathreplacing}
\usepackage[most]{tcolorbox}
\usepackage{enumitem}
\usepackage{booktabs}
\usepackage{xcolor}

\pagestyle{empty}
\setlength{\parindent}{0pt}
\setlength{\parskip}{0.4em}

\definecolor{cel}{HTML}{1F4E79}
\definecolor{regula}{HTML}{E8F1FA}
\definecolor{pulapka}{HTML}{FDECEA}
\definecolor{okres}{HTML}{C0392B}
\definecolor{skonczony}{HTML}{1E8449}

% Ramka na regułę
\newtcolorbox{regula}[1][]{
  colback=regula, colframe=cel, boxrule=0.8pt, arc=2pt,
  left=6pt, right=6pt, top=4pt, bottom=4pt, title=#1,
  fonttitle=\bfseries, coltitle=white
}

% Ramka na pułapkę z kartkówki
\newtcolorbox{pulapka}[1][]{
  colback=pulapka, colframe=okres, boxrule=0.8pt, arc=2pt,
  left=6pt, right=6pt, top=4pt, bottom=4pt, title=#1,
  fonttitle=\bfseries, coltitle=white
}

% Nagłówek karty: numer, tytuł, cel
\newcommand{\naglowekKarty}[3]{%
  {\LARGE\bfseries\color{cel} Karta #1: #2}\\[2pt]
  {\large Cel: #3}\par
  \vspace{0.3em}\noindent\rule{\linewidth}{1pt}\par
}

% Pole na obliczenia ucznia: ramka z delikatną kratką, argument = wysokość
\newcommand{\miejsce}[1]{%
  \par\vspace{0.2em}\noindent
  \begin{tikzpicture}
    \draw[step=0.5cm, gray!25, very thin] (0,0) grid (\textwidth,#1);
    \draw[gray!60, rounded corners=2pt] (0,0) rectangle (\textwidth,#1);
    \node[anchor=north west, font=\scriptsize, text=gray] at (0.05,#1) {obliczenia};
  \end{tikzpicture}\par}
 \begin{document}

% Silnik tectonic (XeTeX): polskie znaki przez fontspec, nie przez inputenc/T1
\usepackage{fontspec}
\setmainfont{Arial}
\usepackage[polish]{babel}
\usepackage[a4paper,margin=1.5cm]{geometry}
\usepackage{amsmath,amssymb}
\usepackage{tikz}
\usetikzlibrary{arrows.meta,positioning,calc,shapes.geometric,decorations.pathreplacing}
\usepackage[most]{tcolorbox}
\usepackage{enumitem}
\usepackage{booktabs}
\usepackage{xcolor}

\pagestyle{empty}
\setlength{\parindent}{0pt}
\setlength{\parskip}{0.4em}

\definecolor{cel}{HTML}{1F4E79}
\definecolor{regula}{HTML}{E8F1FA}
\definecolor{pulapka}{HTML}{FDECEA}
\definecolor{okres}{HTML}{C0392B}
\definecolor{skonczony}{HTML}{1E8449}

% Ramka na regułę
\newtcolorbox{regula}[1][]{
  colback=regula, colframe=cel, boxrule=0.8pt, arc=2pt,
  left=6pt, right=6pt, top=4pt, bottom=4pt, title=#1,
  fonttitle=\bfseries, coltitle=white
}

% Ramka na pułapkę z kartkówki
\newtcolorbox{pulapka}[1][]{
  colback=pulapka, colframe=okres, boxrule=0.8pt, arc=2pt,
  left=6pt, right=6pt, top=4pt, bottom=4pt, title=#1,
  fonttitle=\bfseries, coltitle=white
}

% Nagłówek karty: numer, tytuł, cel
\newcommand{\naglowekKarty}[3]{%
  {\LARGE\bfseries\color{cel} Karta #1: #2}\\[2pt]
  {\large Cel: #3}\par
  \vspace{0.3em}\noindent\rule{\linewidth}{1pt}\par
}

% Pole na obliczenia ucznia: ramka z delikatną kratką, argument = wysokość
\newcommand{\miejsce}[1]{%
  \par\vspace{0.2em}\noindent
  \begin{tikzpicture}
    \draw[step=0.5cm, gray!25, very thin] (0,0) grid (\textwidth,#1);
    \draw[gray!60, rounded corners=2pt] (0,0) rectangle (\textwidth,#1);
    \node[anchor=north west, font=\scriptsize, text=gray] at (0.05,#1) {obliczenia};
  \end{tikzpicture}\par}


\begin{document}

\naglowekKarty{4}{Punkty procentowe czy procenty}{odróżnić „o ile punktów procentowych” od „o ile procent”}

\vspace{0.3em}
\begin{regula}[REGUŁA – DWA RÓŻNE PYTANIA]
\begin{enumerate}[leftmargin=*, itemsep=3pt]
  \item \textbf{O ile punktów procentowych (pp)?} Po prostu odejmij oba procenty. \quad $25\%-20\%=5\ \text{pp}$.
  \item \textbf{O ile procent?} Porównaj zmianę z wartością \emph{początkową}: $\dfrac{\text{zmiana}}{\text{początek}}\cdot 100\%$. \quad $\dfrac{5}{20}=\dfrac{1}{4}=25\%$.
\end{enumerate}
\end{regula}

\vspace{0.4em}
\textbf{Krok po kroku: poparcie wzrosło z $20\%$ do $25\%$.}
\begin{enumerate}[leftmargin=*, itemsep=3pt]
  \item Punkty procentowe: $25-20=5$, więc wzrost o $5$ pp.
  \item Procenty: zmiana $5$ w stosunku do początku $20$: $\dfrac{5}{20}=25\%$, więc wzrost o $25\%$.
\end{enumerate}
\textbf{Druga droga: sto osób.} Gdy zapytano $100$ osób, było $20$ zwolenników, a jest $25$. Przybyło $5$ osób (to $5$ pp). Było ich $20$, a przybyło $5$ z tych $20$, czyli $\dfrac{1}{4}=25\%$.

\begin{center}
\begin{tikzpicture}[>=Stealth]
  % Dwa paski po 100 jednostek (1 jednostka = 0.1 cm)
  \draw[thick] (0,1.4) rectangle (10,2.1);
  \fill[cel!35] (0,1.4) rectangle (2,2.1); \draw[thick] (0,1.4) rectangle (10,2.1);
  \node[font=\small\bfseries, color=cel] at (1,1.75) {$20\%$};
  \draw[thick] (0,0.2) rectangle (10,0.9);
  \fill[cel!35] (0,0.2) rectangle (2,0.9); \fill[okres!40] (2,0.2) rectangle (2.5,0.9); \draw[thick] (0,0.2) rectangle (10,0.9);
  \node[font=\small\bfseries, color=cel] at (1,0.55) {$20\%$};
  \node[font=\small\bfseries, color=okres] at (2.25,-0.15) {$+5\ \text{pp}$};
  \node[font=\small, anchor=west] at (10.2,1.75) {było $20\%$};
  \node[font=\small, anchor=west] at (10.2,0.55) {jest $25\%$};
  \node[font=\small\bfseries, color=okres, anchor=west] at (3.2,0.55) {$5$ z $20$ to $25\%$};
\end{tikzpicture}
\end{center}

\vspace{0.1em}
\textbf{Przykłady rozstrzygnięte:}
\begin{center}
\renewcommand{\arraystretch}{1.7}
\begin{tabular}{@{}lll@{}}
\toprule
Zmiana & O ile pp & O ile \% \\
\midrule
$20\%\to 25\%$ & wzrost o $5$ pp & $\dfrac{5}{20}=25\%$ wzrostu \\[9pt]
$10\%\to 20\%$ & wzrost o $10$ pp & $\dfrac{10}{10}=100\%$ wzrostu (podwoiło się) \\[9pt]
$50\%\to 40\%$ & spadek o $10$ pp & $\dfrac{10}{50}=20\%$ spadku \\[9pt]
$5\%\to 4\%$ & spadek o $1$ pp & $\dfrac{1}{5}=20\%$ spadku \\
\bottomrule
\end{tabular}
\end{center}

\begin{pulapka}[PUŁAPKI Z KARTKÓWKI]
\begin{enumerate}[leftmargin=*, itemsep=2pt]
  \item „Wzrost o $5\%$” i „wzrost o $5$ pp” to dwie różne rzeczy. Zdanie „z $20\%$ na $25\%$, czyli o $5\%$” jest fałszywe.
  \item Procenty (zmiana względna) dzielimy zawsze przez wartość \emph{początkową}, nie końcową.
  \item Punkty procentowe odejmujemy, bez żadnego mnożenia. Mnożenie pojawia się dopiero przy „o ile procent”.
\end{enumerate}
\end{pulapka}

\end{document}
